\section{Hierarchical decomposition}
\label{sec:hierarchical}

% ---------------------------------------------------------------
\subsection{Task-space / null-space at each phase}
\label{sec:hierarchical-tn}

\Cref{def:TS-NS} of \cref{sec:task-null} introduced the
task-space subspace
$\TS = \mathrm{range}(J^{\top}(\jointvec))$ and the null-space
subspace $\NS = \ker(J(\jointvec))$ as a configuration-dependent
splitting of the joint-velocity space $\mathbb{R}^{n}$.
\Cref{sec:phase-cycle} established the five-phase decomposition
of a pick--place trajectory together with the structural
operators $\phaseop{0}, \ldots, \phaseop{4}$ associated with the
phases. The present subsection ties the two together. For each
phase $k \in \{1, \ldots, 5\}$ we identify the \emph{task
variable}---what the manipulator is required to control on $\TS$
during that phase---and the \emph{null-space freedom}---what the
manipulator is free to vary in $\NS$ without affecting the task
variable. The per-phase reading is the structural prerequisite
for the bi-level specialization of \cref{sec:hierarchical-bilevel}
and for the dexterity commutator of
\cref{sec:hierarchical-commutator}.

\paragraph{Per-phase specialization.}
The task variable on $\TS$ and the null-space freedom on $\NS$
specialize phase-by-phase as in \cref{tab:tn-per-phase}. Phases 1
and 5 are out-of-contact reaches, so the task variable is the
end-effector pose alone; phases 2 and 4 are the bang arcs at
$t_{1}$ and $t_{4}$, so the task variable is the contact wrench
at the corresponding $\swsurf$-crossing; phase 3 is the in-contact
transport, so the task variable is the object pose trajectory.

\begin{table}[h]
\centering
\small
\begin{tabular}{@{}clll@{}}
\toprule
$k$ & Phase & Task variable on $\TS$ & Null-space freedom on $\NS$ \\
\midrule
1 & pick (approach) & end-effector pose $g \in SE(3)$
  & redundant arm configuration (e.g.\ elbow swivel) \\
2 & load (contact establishment) & contact-wrench at $\rcontact = 1^{+}$
  & wrist orientation consistent with grasp geometry \\
3 & use (transport) & object pose trajectory $x_{\mathrm{obj}}(t) \in SE(3)$
  & redundant gripper-arm trajectory \\
4 & unload (release) & contact-wrench at $\rcontact = 1^{-}$
  & wrist orientation as in phase 2 \\
5 & place (retreat) & end-effector pose $g \in SE(3)$
  & redundant arm configuration as in phase 1 \\
\bottomrule
\end{tabular}
\caption{Per-phase specialization of the task-space /
null-space split of \cref{def:TS-NS}. Phases 1 and 5 control the
end-effector pose; phases 2 and 4 control the contact wrench at
the $\swsurf$-crossings; phase 3 controls the object pose
trajectory through the contact.}
\label{tab:tn-per-phase}
\end{table}

\paragraph{Per-phase dimension count for the Franka Panda.}
For the concrete instance used in the experimental staircase of
\cref{sec:experiments}, the manipulator is the Franka Emika Panda
with $n = 7$ joints. In phases 1, 3, and 5 the task variable is a
full $SE(3)$ pose (end-effector pose, object pose, end-effector
pose, respectively), so $m = 6$ and the dimensions of the
splitting are
\[
  \dim \TS \;=\; 6,
  \qquad
  \dim \NS \;=\; n - m \;=\; 1,
\]
with the single null-space direction the elbow-swivel self-motion
of \cref{ex:franka-elbow}. In phases 2 and 4 the task variable
is the contact wrench at the $\swsurf$-crossing; in the
three-dimensional world the wrench has six independent components
(three force, three moment), so the task dimension is again
$m = 6$ and $\dim \NS = 1$, which is now realized by the wrist
orientation freedom that maintains the grasp normal aligned with
the contact face. (For a planar grasp the wrench dimension drops
to three and $\NS$ correspondingly enlarges; the Franka platform
with a parallel-jaw end-effector executes the planar specialization
when the manipulated face is vertical.) The same arm therefore
exposes the same one-dimensional self-motion in every phase, but
the physical interpretation of that self-motion changes with the
phase: an elbow swivel in phases 1, 3, 5 and a wrist-orientation
freedom in phases 2 and 4.

\paragraph{Phase 3 in detail.}
Phase 3 is the substantive in-contact transport phase and the
only phase where the task variable is the object pose rather than
the end-effector pose. The manipulator must track an object pose
trajectory $x_{\mathrm{obj}}(t) \in SE(3)$ on the open interval
$(t_{1}, t_{4})$ through the bi-level coupling of
\cref{def:bilevel}: the outer level prescribes
$x_{\mathrm{obj}}(t)$ and the inner level chooses contact forces
that realize the corresponding object acceleration while
maintaining $\specgap(t) \ge \varepsilon$. The null-space freedom
on $\NS$ is the redundant gripper-arm trajectory: multiple arm
configurations realize the same object pose, and the controller
exploits this redundancy to satisfy the spectral-gap constraint
$\specgap \ge \varepsilon$ of \cref{eq:bilevel-outer} while
remaining strictly inside the friction cone $\frictioncone$. This
is the per-phase realization of the operational-space framework
of \citet{Khatib1987} and \citet{Nakamura1991}, specialized to
the contact-mediated transport regime: the same orthogonal
projector $P_{\NS} = I - J^{+} J$ that
\eqref{eq:qdot-split} of \cref{sec:task-null} writes generally is
applied here to the contact-coupled dynamics of phase 3, with the
null-space command directed by the inner force-closure program.

\paragraph{Forward references.}
\Cref{sec:hierarchical-bilevel} specializes the abstract bi-level
scaffold of \cref{def:bilevel} and the coupling proposition
\cref{prop:bilevel-coupling} to the per-phase task and null
variables identified above, deriving the explicit form of the
acceleration bound $\alpha(\jointvec, f_{n,\max})$ for the
parallel-jaw grasp under wind. \Cref{sec:hierarchical-commutator}
defines the commutator $[\TS, \NS]_{\jointvec}$ that quantifies
the dynamic dexterity of the manipulator on its null-space and
proves that this commutator is non-zero exactly in the regime
where internal reconfiguration modulates task-space forces.

% ---------------------------------------------------------------
\subsection{Bi-level structure for the pick--place cycle}
\label{sec:hierarchical-bilevel}

\Cref{def:bilevel} of \cref{sec:bilevel} stated the abstract
bi-level scaffold for the pick--place OCP: an outer
trajectory-optimization level over the object pose
$x_{\mathrm{obj}}(t)$ and an inner force-closure level over
contact forces $f(t)$ at every instant. The coupling
proposition \cref{prop:bilevel-coupling} stated, in addition,
that force-closure feasibility induces a state-dependent
acceleration bound
$\|\ddot{x}\| \le \alpha(\jointvec, f_{n,\max})$ on the outer
program but left the explicit form of the function $\alpha$
unspecified. The present subsection derives that explicit form
for the parallel-jaw grasp under wind, which is the worked
instance the experimental staircase of \cref{sec:experiments}
will exercise. The deliverable is the explicit
\cref{prop:accel-bound-explicit} below; the per-phase reading
of the bi-level coupling closes the subsection.

\paragraph{Setup of the parallel-jaw grasp under wind.}
The two-finger parallel-jaw gripper of
\cref{ex:parallel-jaw} of \cref{sec:rho} is treated here under
an additional adversarial horizontal wind force $\Fwind(t)$
acting on the manipulated body, which is the perturbation regime
of \cref{sec:curriculum}. Each finger applies a normal force
of magnitude $f_{n,L}$ and $f_{n,R}$, respectively, to opposite
vertical faces of a rigid box of mass $m$ in a uniform
gravitational field of magnitude $\|g\|$. The grasp is symmetric:
$f_{n,L} = f_{n,R} = f_{n}$. The friction coefficient at each
finger contact is $\mu$. The static-equilibrium normal-force
threshold of \cref{def:rho} is now sourced by gravity together
with the wind: the box is in static equilibrium when the total
tangential friction available at the two contacts dominates
the combined load $m \|g\| + \|\Fwind(t)\|$. Reproducing the
parallel-jaw order-parameter of \cref{eq:rho-grasp} with the
wind term added,
\begin{equation}
\label{eq:rho-grasp-wind}
  \rcontact_{\mathrm{grasp}}(t)
  \;=\;
  \frac{2\, \mu\, f_{n}(t)}{m \|g\| + \|\Fwind(t)\|},
  \qquad
  \rcontact_{\mathrm{grasp}}^{\star} \;=\; 1,
\end{equation}
which is the natural extension of \cref{ex:parallel-jaw} to the
wind-perturbed regime. The denominator is now time-varying
through $\Fwind(t)$, and equality to one is the slipping
threshold of the wind-perturbed grasp.

\paragraph{Inner level: maximum spectral gap admissible under
the friction cone.}
The contact graph of the symmetric two-finger grasp is
$G_{c} = (\{b, L, R\},\, \{(b, L), (b, R)\})$ with edge weights
$w_{bL} = w_{bR} = w$ given by the contact stiffness, which is
proportional to the normal force $f_{n}$ in the linearized
contact model of \cref{sec:specgap}. Ordering the vertices
$b, L, R$, the contact-graph Laplacian
\eqref{eq:laplacian-def} reads
\begin{equation}
\label{eq:two-finger-L}
  L_{c} \;=\;
  \begin{pmatrix}
    2 w & -w & -w \\
    -w & w & 0 \\
    -w & 0 & w
  \end{pmatrix},
\end{equation}
whose spectrum is $\{0,\, w,\, 3 w\}$. The spectral gap
\eqref{eq:specgap-def} is therefore
$\specgap(L_{c}) = w$, equal to the contact stiffness at each
finger. Substituting back into the inner program
\eqref{eq:bilevel-inner}, the inner level chooses $f_{n}$ to
maximize $\specgap = w(f_{n})$ subject to the friction-cone
constraint $|f_{t}| \le \mu f_{n}$ and the Newton balance
\eqref{eq:newton-balance}. Given the box mass $m$, gravity
$g$, wind $\Fwind$, and the desired object acceleration
$\ddot{x}_{\mathrm{des}}$ all prescribed by the outer level, the
Newton balance pins the resultant contact force uniquely; the
inner level reduces to checking that this pinned force lies
inside the friction cone, and the maximum-spectral-gap solution
is the inner solution that realizes the pinned $f_{n}$ value.
This is the operational form of force-closure analysis in the
sense of \citet{MurrayLiSastry1994}, specialized to the
two-finger grasp.

\paragraph{Outer level: derivation of the explicit
acceleration bound.}
The substantive deliverable of the present subsection is the
explicit form of the acceleration bound
$\alpha(\jointvec, f_{n,\max})$ promised by
\cref{prop:bilevel-coupling}~(ii) of \cref{sec:bilevel}.
Project the Newton balance \eqref{eq:newton-balance} onto the
direction tangent to the contact normal at each finger, and let
$\|\cdot\|_{t}$ denote the tangential component along that
direction. The tangential projection of the resultant force
required to accelerate the box at $\ddot{x}_{\mathrm{des}}$
against the wind is
\begin{equation}
\label{eq:ft-projection}
  |f_{t}|
  \;=\;
  m \,\|\ddot{x}_{\mathrm{des},\, t}\|
  \;+\; \|\Fwind(t)\|_{t},
\end{equation}
where the first term is the inertial contribution and the
second the wind contribution along the tangential direction.
The friction-cone constraint $|f_{t}| \le \mu f_{n}$ together
with the per-finger normal-force budget $f_{n} \le f_{n,\max}$
(a hardware ceiling on the gripper actuator) then bounds the
tangential acceleration:
\begin{equation}
\label{eq:accel-bound-derivation}
  m \,\|\ddot{x}_{\mathrm{des},\, t}\|
  \;\le\;
  \mu \, f_{n,\max} \;-\; \|\Fwind(t)\|_{t}.
\end{equation}
Dividing by $m$ yields the explicit form of $\alpha$
promised in \eqref{eq:accel-bound}. The dependence on the
configuration $\jointvec$ enters through the orientation of the
contact normal in the world frame, which determines how much of
the wind force projects into the tangential direction.

\begin{proposition}[Acceleration bound for the parallel-jaw
grasp under wind]\label{prop:accel-bound-explicit}
For a symmetric parallel-jaw grasp of a rigid box of mass $m$
with friction coefficient $\mu$ at each finger, normal-force
budget $f_{n,\max}$ per finger, and wind force $\Fwind(t)$, the
acceleration bound \eqref{eq:accel-bound} of
\cref{prop:bilevel-coupling} takes the explicit form
\begin{equation}
\label{eq:alpha-explicit}
  \alpha(\jointvec, f_{n,\max})
  \;=\;
  \frac{\mu\, f_{n,\max} \;-\; \|\Fwind(t)\|_{t}(\jointvec)}{m},
\end{equation}
where $\|\cdot\|_{t}$ denotes the tangential component along the
gripper contact normal at configuration $\jointvec$.
\end{proposition}

\begin{proof}
The friction-cone constraint $|f_{t}| \le \mu f_{n}$ together
with the per-finger normal-force ceiling $f_{n} \le f_{n,\max}$
gives $|f_{t}| \le \mu f_{n,\max}$. Project the Newton balance
\eqref{eq:newton-balance} of \cref{sec:bilevel} onto the
direction tangent to the contact normal: the tangential
component of the resultant required to accelerate the box
against the wind is given by \eqref{eq:ft-projection}.
Substituting into the cone bound,
$m \|\ddot{x}_{\mathrm{des},\, t}\| + \|\Fwind\|_{t} \le \mu
f_{n,\max}$, which rearranges to \eqref{eq:alpha-explicit}.
The configuration dependence of the right-hand side is through
the wind-tangential component
$\|\Fwind(t)\|_{t}(\jointvec)$, which depends on the orientation
of the contact normal at $\jointvec$.
\end{proof}

\paragraph{Per-phase reading of the bi-level coupling.}
The bi-level coupling derived above is non-trivial only on
phase 3: phases 1 and 5 are out-of-contact reaches with no
robot--object contact, so the inner program is not posed and
the outer trajectory is constrained only by the manipulator
kinematics. Phase 3 is the in-contact transport on the open
interval $(t_{1}, t_{4})$, and the outer trajectory must
respect the inner force-closure feasibility throughout: the
acceleration bound \eqref{eq:alpha-explicit} is active for
every $t \in (t_{1}, t_{4})$, and the spectral-gap maintenance
condition \eqref{eq:bilevel-outer} likewise holds on the
in-contact interior. Phases 2 and 4 are the bang arcs at
$t_{1}$ and $t_{4}$ at which the inner program transitions
across its singular boundary: at $t_{1}$ the inner solution
crosses from $\specgap = 0$ up to $\specgap = \varepsilon$ and
at $t_{4}$ it crosses back down. The bang force $u_{s}$ of
\cref{def:three-term} is exactly the impulsive contribution
that drives $\specgap$ across the threshold $\varepsilon$ at
each crossing. The explicit acceleration bound
\eqref{eq:alpha-explicit} is therefore the constraint on the
outer program throughout phase 3, whereas phases 2 and 4 are
governed by the bang structure of \cref{thm:three-term} rather
than by a continuous acceleration bound.
